\chapter{Binomial Coefficients}

\section{Combinatorial Proof}

\textbf{Set-up}
\begin{itemize}
    \item Given an equation to prove. $LHS=RHS$
    \item Come up with some set of objects one of the sides (say $LHS$) is equivalent to in size.
    \item Come up with an alternative way to count the collection of objects that corresponds to the other side($RHS$).
    \item Conclude that you have counted the same set in two different ways to deduce $LHS=RHS$
\end{itemize}
Basically, it is proof by story telling.

\begin{theorem}
    $\forall n \in \N. 2^n=\sum_{k=0}^{n}{n \choose k}$
\end{theorem}
\begin{proof}
(The first combinatorial proof)
    \begin{itemize}
        \item Consider the collection of length $n$ binary strings. There are $2^n$ many such strings.
        \item Note, we can seperate strings in accordance with how many 1's the have. For $k \in \{0, \cdots ,n\}$, there are exactly ${n \choose k}$ many strings of length $n$ with k 1's. Therefore, there are $\sum_{k=0}^{n}{n \choose k}$ many binary strings.
        \item Thus, $2^n=\sum_{k=0}^{n}{n \choose k}$ as desired.
    \end{itemize}
\end{proof}
\begin{proof}
(The second combinatorial proof)
\begin{itemize}
    \item Suppose I am having a party and can invite any (or none) of my $n$ friends. There are $2^n$ many ways to do this, as there are $2^n$ many bsubsets from a collection of $n$ people.(for each people, choose him or not)
    \item Note, I can invite any number from $k=0, \cdots ,n$ of my friends. There are $2^n=\sum_{k=0}^{n}{n \choose k}$ many ways to choose k friends from n. Therefore, there are $\sum_{k=0}^{n}{n \choose k}$ many ways.
    \item Thus, $2^n=\sum_{k=0}^{n}{n \choose k}$ as desired.
\end{itemize}
\end{proof}

\section{Stars and Bars}
\begin{theorem}
    Let n, k be positive integers. There are ${n+(k-1) \choose k-1}$ integer solutions to the equatioin 
    $x_1 + \cdots + x_k=n$ with constraint $x_i \geq 0$ for all i.
\end{theorem}
\begin{proof}
    (This is not rigorous, I can only give you some intuition)\\
    Suppose there are n objects (represented by stars) to be placed into k bins. 
    The bins are distinguishable (say they are numbered 1 to k) but the n stars are not (so configurations are only distinguished by the number of stars present in each bin). 
    A configuration is thus represented by a k-tuple of non-negative integers, as in the statement of the theorem. Start by placing the stars in a line,
    we can place multiple bars between stars, before the first star and after the last star. 
    Since that any arrangement of stars and bars consists of a total of $n + k - 1$ objects, n of which are stars and $k-1$ of which are bars.
    We only need to choose $k - 1$ of the $n + k - 1$ positions to be bars (or, equivalently, choose n of the positions to be stars).
    Thus, there are ${n+(k-1) \choose k-1}$ arrangements.
\end{proof}

\begin{corollary}
    Let n, k be positive integers. There are ${n-1 \choose k-1}$ integer solutions to the equatioin 
    $x_1 + \cdots + x_k=n$ with constraint $x_i > 0$ for all i.
\end{corollary}
\begin{proof}
    The question is equivalent to find the number of solutions of $y_1 + \cdots + y_k=n-k$ with constraint $y_i \geq 0$ for all i.
    Then use the theorem above finish the proof.
\end{proof}

\section{Binomial Theorem}

\begin{theorem}
    $\forall x \in \R. (1+x)^n=\sum_{k=0}^{n}{n \choose k}x^k$
\end{theorem}

\begin{proof}
    We know that the function $f(x)=(1+x)^n$ is a polynomial, 
    so $f(x)=\sum_{k=0}^{n}c_kx^k$ for some $c_k \in \Z$. 
    Consider the derivative of $f(x)=(1+x)^n$, then the i-th derivative is $$f^{(i)}(x)=P(n,i)(1+x)^{n-i}$$ 
    Computed another way, we have $$f^{(i)}(x)=\sum_{k=1}^{n}c_kP(k,i)x^{k-i}$$
    Evaluating $f^{(i)}(0)$, we have $P(n,i)=c_iP(i,i)$. 
    Hence we get $$c_i=\frac{P(n,i)}{P(i,i)}=\frac{n!}{(n-i)!}\frac{(i-i)!}{i!}={n \choose i}$$ as desired.
\end{proof}

\begin{theorem}[Full Binomial Theorem]
    $\forall x \in \R. (x+y)^n=\sum_{k=0}^{n}{n \choose k}x^ky^{n-k}$
\end{theorem}
\begin{proof}
    It is true when $y=0$, so suppose $y \neq 0$. 
    Note $(x+y)^n=y^n(1+\frac{x}{y})^n$ ,so set $z=\frac{x}{y}$. 
    finally we have 
    \begin{equation}
        \begin{split}
            y^n(1+\frac{x}{y}) &= y^n \sum_{k=0}^{n}{n \choose k}z^k\\
            &= y^n \sum_{k=0}^{n}{n \choose k}x^ky^{-k}\\
            &= \sum_{k=0}^n{n \choose k}x^ky^{n-k}
        \end{split}
    \end{equation}
\end{proof}

\section{Lattice and Catalan number}

\begin{definition}
    A \textbf{lattice} is a set of the form $[n+1] \times [m+1]$. 
    We call a lattice of this form \textbf{an n by m lattice}. 
    The point $(0,0) \in [n+1]\times [m+1]$ is called the  \textbf{origin}.
\end{definition}

\begin{definition}
    A \textbf{lattice path} is a swquence $\{(a_0,b_0),(a_1,b_1), \cdots ,(a_{n+m},b_{n+m})\}$ 
    that satisfies the property that:\\
    \begin{itemize}
        \item $(a_0,b_0)=(0,0)$
        \item $(a_{n+m},b_{n+m})=(n,m)$
        \item $(m_{i+1},n_{i+1})=(m_{i+1},n_i)$ or $(m_{i+1},n_{i+1})=(m_{i},n_{i+1})$
    \end{itemize}
\end{definition}

\begin{theorem}
    Given an n by m lattice, there are ${n+m \choose n}$ many ways to walk from the origin to $(n,m)$.
\end{theorem}
\begin{proof}\
    \begin{itemize}
        \item Nottice that a lattice path is codified by an length $n+m$ sequence from $\{U,R\}$ with n many R's(U means up and R means right)
        \item To see this, notice the following reversible process to get a walk from a sequence and vice versa.
        \item Given a lattice path $P=\{(a_o.b_0),(a_1,b_1), \cdots ,(a_{n+m},b_{n+m})\}$, assign to it a sequence $s:[n+m] \rightarrow \{U,R\}$ where $s(i)=U$ iff $b_{i+1}=b_i+1$.
        \item To reverse this, given a sequence $s:[n+m] \rightarrow \{U,R\}$, assign to it a lattice path $P=\{(a_o.b_0),(a_1,b_1), \cdots ,(a_{n+m},b_{n+m})\}$ where $b_{i+1}=b_i+1$ when $s(i)=U$ or alternatively, $a_{i+1}=a_i+1$ when $s(i)=R$
    \end{itemize}
\end{proof}

\begin{definition}
    We define the \textbf{Catalan number} $C(n)$ to be the number of lattice paths from $(0,0)$ to $(n,n)$ that stay below the line $y=x$.
\end{definition}

\begin{theorem}
    The catalan number can be explicitly computed,
    $$C(n)=\frac{1}{n+1}{2n \choose n}$$
\end{theorem}
\begin{proof}
(proof1, Total - Above Diagonal)
\begin{itemize}
    \item Total= ${2n \choose n}= \frac{(2n)!}{n!n!}$
    \item Above= ${2n \choose n+1}= \frac{(2n)!}{(n+1)!(n-1)!}$ (this part is very tricky, I will not give details)
    \item $C(n)={2n \choose n}-{2n \choose n+1}=\frac{1}{n+1}{2n \choose n}$
\end{itemize}
\end{proof}
\begin{proof}
(proof2, Find the recurrence relation)\\
Let $k \in \{1, \cdots ,n\}$ be arbitrary, consider the spot $(k,k)$.\\
Let $S_k$ denotes the set of paths which first intersect the line $y=x$ at point $(k,k)$.\\
For $k=0$ the only path is stay at the origin, hence $c_0=|S_0|=1$.\\
For $k \in \{1, \cdots ,n\}$, can check that $|S_k|=c_{k-1} \cdot c_{n+1-k}$\\
Since $\{S_0, \cdots ,S_k\}$ forms a partition of $S_{n+1}$, hence we get the following recurrence
\begin{equation}
    \begin{cases}
        c_0=1\\
        c_{n+1}=\sum_{i=0}^{n}c_ic_{n-i}
    \end{cases}
\end{equation}
We can check the example of \ref{cnr} to finish the proof.
\end{proof}
\noindent Example: How many ways can we parenthesize the string abcdef?